\section{Independent review findings}
\label{sec:review}

An independent reviewer, who did not produce the work, reviewed revision
\texttt{72f67f0} (task T-060) \ev{reports/review/72f67f0/findings.md} and then
re-reviewed the fixes in \texttt{72f67f0..05a06a3} (task T-061)
\ev{reports/review/05a06a3/findings.md}. Dispositions are recorded in ADR-014 and
ADR-018 \ev{management/DECISIONS.md}.

\paragraph{Summary.}
\begin{itemize}[noitemsep]
  \item First review: 16 findings (0 critical, 3 major, 9 minor, 4 info);
        dispositions 13 FIX, 1 ACCEPTED, 2 info FIX (ADR-014).
  \item Re-review: 14 CLOSED, 1 ACCEPTED (F-012), 1 NOT CLOSED at 05a06a3 (F-007,
        fixed afterwards, ADR-018); 5 new findings (2 minor, 3 info).
  \item Reviewer re-checks at 05a06a3: \texttt{run\_pytest tests/firmware}
        ``295 passed in 1.93s''; ERC ``Found 0 violations''; DRC 0/0/0; status
        counts 17/18/3/4 = 42 match (\lbl{MEASUREMENT}).
\end{itemize}

\paragraph{Highlight: defect found while fixing F-005.} On an I\textsuperscript{2}C
NACK the Adafruit library's \texttt{touched()} returned 0x0FFF, i.e.\ all 12 pads
touched, so a lost sensor would have sent 12 Note Ons. The firmware now reads the
status registers itself and checks the transaction result (ADR-017). The re-review
confirmed the fix (F-005 CLOSED).

{\footnotesize
\begin{longtable}{@{}>{\raggedright\arraybackslash}p{0.07\linewidth}>{\raggedright\arraybackslash}p{0.07\linewidth}>{\raggedright\arraybackslash}p{0.40\linewidth}>{\raggedright\arraybackslash}p{0.23\linewidth}>{\raggedright\arraybackslash}p{0.17\linewidth}@{}}
\caption{Review findings, dispositions and re-review verdicts.}\label{tab:review}\\
\toprule
ID & Sev. & Finding (one line) & Disposition & Re-review \\
\midrule
\endfirsthead
\toprule
ID & Sev. & Finding (one line) & Disposition & Re-review \\
\midrule
\endhead
\bottomrule
\endlastfoot
F-001 & major & A wrong JP1 position was called harmless; it exceeds the MPR121 VDD + 0.3\,V absolute maximum and back-feeds the module. & FIX: README \S5/\S6 & CLOSED \\
F-002 & major & The $\pm$4\,kV ``PASS'' compared peak current only; pin limit is an \lbl{ASSUMPTION}; resistor not pulse-rated. & FIX: REQ-023 is an estimate, not a pass & CLOSED \\
F-010 & major & REQ-018 was PARTIAL although budget check A3b FAILs (4.40\,V $<$ 4.50\,V). & FIX: REQ-018 FAIL & CLOSED \\
F-003 & minor & Pot hysteresis defined two ways; offset share misstated. & FIX: REQ-006 restated & CLOSED \\
F-004 & minor & Two held pads can share a note; releasing one sent Note Off. & FIX: reference counting & CLOSED \\
F-005 & minor & Sensor lost during play not detected; missing second sensor never retried. & FIX: sensor health & CLOSED \\
F-006 & minor & Leonardo could run I\textsuperscript{2}C at 100\,kHz after \texttt{begin()}. & FIX: \texttt{setClock} after every \texttt{begin()} & CLOSED \\
F-007 & minor & REQ-033 build logs not in the commit. & FIX: commit logs & NOT CLOSED at 05a06a3; fixed after (ADR-018) \\
F-008 & minor & budgets.md stale; coverage.md ``USL 201''. & FIX: regenerate, reconcile & CLOSED (residue N-002) \\
F-009 & minor & README named an ``aprender'' button. & FIX & CLOSED \\
F-011 & minor & Latency model omitted the 15\,ms pad lockout. & FIX: REQ-014 scoped to first touch & CLOSED \\
F-014 & minor & Terminal pin diagonal 1.27\,mm in a 1.3\,mm hole; pot slot drilling unclear. & FIX: 1.4\,mm drill, slot instruction & CLOSED \\
F-012 & info & Latency rests on an assumed MPR121 status-update model. & ACCEPTED (plan P-2) & ACCEPTED \\
F-013 & info & ERC ignored single-pin labels globally. & FIX at root (J3) & CLOSED \\
F-015 & info & BOM lacked part numbers; strip count wrong. & FIX & CLOSED \\
F-016 & info & AGND relied on the Pico's internal tie without citation. & FIX: cite S-01/S-02 & CLOSED \\
N-001 & minor & Automatic restart of a self-reset sensor blocks the loop 16.8--22.1\,ms, beyond REQ-016's 10\,ms. & ACCEPTED as specified exception (ADR-018) & new \\
N-002 & minor & REQ-016 evidence stale (worst loop $\approx$1.72\,ms incl.\ ECR poll). & FIXED (ADR-018) & new \\
N-003 & info & INTERFACES had no J3/J4--J7. & FIXED (frozen-3, IF-PCB-AUX) & new \\
N-004 & info & With 2 sensors configured and one absent the LED blinks fast permanently. & ACCEPTED (matches REQ-009) & new \\
N-005 & info & A second pad on an already sounding note is silent (no retrigger). & ACCEPTED (matches REQ-042) & new \\
\end{longtable}
}

\paragraph{What the review did not cover.} First review: INTERFACES.json field by
field; visual check of the copper PDFs and mirroring; the sampler beyond its MIDI
handler; 8 of 9 firmware targets; MPR121 power-on-reset timing. Re-review: compile
runs (logs only); copper PDFs and silkscreen visually; this report's LaTeX; the
individual test cases of \texttt{test\_notes\_health.py}. Neither review is a
hardware test.

\subsection{CR-002 independent re-review}
\label{sec:cr002-review}
An independent reviewer re-reviewed the CR-002 outputs at \texttt{8394b41} on branch
\texttt{change/CR-002} (task T-071): R1--R12 part and footprint, relayout, gap
check, pulse-rating calculation, budgets revision~6, REQUIREMENTS draft-6,
INTERFACES frozen-4 and the documentation. The harness refused the reviewer's write
of the findings file; the manager recorded the findings from the reviewer's
hand-back report without changing their substance
\ev{reports/review/CR-002/findings.md}. The electronics fixes were made in T-073
(committed in c901ad8) \ev{management/handoffs/electronics-engineer/T-073.md};
dispositions are recorded in ADR-020 \ev{management/DECISIONS.md}.

\paragraph{Summary.}
\begin{itemize}[noitemsep]
  \item 13 findings: 0 critical, 0 major, 6 minor, 7 info.
  \item Dispositions: 8 fixed (C2-006 by systems in T-074, not yet in c901ad8),
        1 partly fixed with a residual accepted risk (C2-002), 1 accepted risk
        (C2-005), 2 accepted (C2-012 with a test added, C2-013), 1 no action
        (C2-004).
  \item Reviewer re-checks at 8394b41, all PASS (recorded in the findings file;
        raw outputs in \nolinkurl{reports/review/CR-002/}): ERC 0; DRC 0,
        unconnected 0, parity 0; \nolinkurl{check_gap.py} PASS (copper only);
        \nolinkurl{calc_flashover} and \nolinkurl{calc_pulse_rated} outputs
        byte-identical; interface check 80/80 (committed output stale, 79/80);
        \nolinkurl{pcb_rules_check} 16/16; budgets byte-identical
        (23.78--37.72\,EUR); independent chart readings at 3.3\,k$\Omega$ 4086 /
        1740 / 1469\,V for 0414 / 0207 / 0204; drill count 131 equal to the
        assembly guide; configuration wires 21.1--32.1\,mm (\lbl{MEASUREMENT},
        design file).
  \item Reviewer's requirement cross-check: REQ-023 PARTIAL (agreed); REQ-027
        PASS; REQ-028/029 PASS with stale evidence text (C2-006); REQ-032 FAIL;
        REQ-036 should be PARTIAL until C2-001 and C2-007 are fixed (both fixed
        in T-073).
\end{itemize}

\paragraph{Highlights.} The closest metal of a flat resistor is its end caps, not
the copper (C2-002): the 8\,kV gap claims now apply to copper only, and the 4\,kV
end-cap result is PARTIAL on an assumed cap length (Section~\ref{sec:esd}). The
pulse energy had been overstated about twice (C2-003): 182\,mJ, not 364\,mJ, with
the verdicts unchanged. A metal screw at the old H3 position could bypass R12
(1.94\,mm bypass sum, FAIL); H3 moved to (96, 38)\,mm (15.36\,mm, PASS)
(\lbl{CALCULATION}, \nolinkurl{check_gap.out}).

{\footnotesize
\begin{longtable}{@{}>{\raggedright\arraybackslash}p{0.07\linewidth}>{\raggedright\arraybackslash}p{0.06\linewidth}>{\raggedright\arraybackslash}p{0.45\linewidth}>{\raggedright\arraybackslash}p{0.36\linewidth}@{}}
\caption{CR-002 re-review findings and dispositions (ADR-020).}\label{tab:review-cr002}\\
\toprule
ID & Sev. & Finding (one line) & Disposition \\
\midrule
\endfirsthead
\toprule
ID & Sev. & Finding (one line) & Disposition \\
\midrule
\endhead
\bottomrule
\endlastfoot
C2-001 & minor & Assembly guide asked for a lead bend $\ge$1.5\,mm from the body; at 15.24\,mm pitch the lead centre is 1.67\,mm from the body end (\lbl{CALCULATION}). & Fixed: reworded (Vishay $M_\mathrm{min}$ 15.0\,mm $\le$ 15.24\,mm pitch; hold the lead next to the body, test-fit one part). \\
C2-002 & minor & Closest metal of a flat resistor is cap to cap, not the 13\,mm copper gap; \nolinkurl{check_gap.py} measures copper only; 8\,kV gap claims overreach. & Partly fixed, residual accepted risk: cap-to-cap check added; 4\,kV PASS for placement error $\le\pm$0.25\,mm (worst 3.61\,mm), FAIL corner 3.13\,mm at $\pm$0.5\,mm; 8\,kV caps FAIL; leads formed on a gauge, caliper check in the plan; 8\,kV statements limited to copper. \\
C2-003 & minor & Pulse energy overstated about 2$\times$ ($\int V\,dt$ form); rating energy-limited at 3.3\,k$\Omega$. & Fixed: 182\,mJ ($\int V^2\,dt$); energy-limited below the 5073\,$\Omega$ knee, plateau 5054\,V; verdicts unchanged. \\
C2-004 & info & Independent chart reading confirms 4086 / 1740 / 1469\,V; dimensions, LSR37 data and part number agree. & No action. \\
C2-005 & minor & Every E terminal pad is 1.40\,mm from the GND edge track, so a strike most likely arcs to GND; the current through the GND trunk to Pico, module and USB was not analysed. & Accepted risk: intentional spark gap; ground bounce 87--339\,V for $\approx$1\,ns, $\le$1.53\,A and $\le$0.92\,nC per 220\,$\Omega$ I\textsuperscript{2}C resistor (\lbl{CALCULATION} on \lbl{ASSUMPTION}ed inductances); arc to GND added to the ESD test; GP26/GP27 HBM rating not confirmed (open). \\
C2-006 & minor & Stale evidence: committed \nolinkurl{interface_check.out} 79/80; REQ-028/029, requirements.csv and INTERFACES frozen-4 notes describe 3a3684b. & Fixed by systems (T-074): \nolinkurl{sim/systems/*.out}, REQUIREMENTS, INTERFACES frozen-5 (committed after c901ad8). \\
C2-007 & minor & Assembly guide offered any resistor lead off-cut for JP1 (0.8\,mm holes); R1--R12 leads are 0.8\,mm. & Fixed: solder or an off-cut of an R13--R16 lead (0.5\,mm) only. \\
C2-008 & info & 1.0\,mm drill for a 0.8\,mm lead sits exactly on the 0.2\,mm fit threshold; generator comment wrong. & Fixed: R1--R12 drill 1.1\,mm (ring 0.55\,mm); comment corrected. \\
C2-009 & info & Courtyard gaps 0.06\,mm (R row to J2) and 0.09\,mm (R12 to H3); M3 head about 0.5--0.9\,mm from R12. & Fixed: H3 moved (95, 48) $\rightarrow$ (96, 38)\,mm; bypass 1.94\,mm FAIL $\rightarrow$ 15.36\,mm PASS; nylon hardware recommended. \\
C2-010 & info & Terminal labels ``10'' and ``11'' touch R10/R11 pad 2 (about 0.1\,mm). & Fixed: worst silkscreen-to-pad gap 0.43\,mm. \\
C2-011 & info & \nolinkurl{check_gap.py} worst-pair name not deterministic on ties. & Fixed: tie-break by reference, then pad. \\
C2-012 & info & I\textsuperscript{2}C relayout under the module (SDA 3.2\,mm from electrode pads over about 31\,mm) not analysed. & Accepted, test added: I\textsuperscript{2}C traffic against baseline noise (V-27). \\
C2-013 & info & IEC 60664-1 F.2 values recalled; the 3.5\,mm rule comes from model M4 ($\eta$ = 0.3). & Accepted: F.2 kept as a cross-check marked ``recalled'' (EE-S25). \\
\end{longtable}
}

\paragraph{Re-checks after the fixes (T-073, c901ad8).} ERC ``ERC messages: 0
Errors 0 Warnings 0''; DRC ``Found 0 DRC violations'', ``Found 0 unconnected
pads''; ``SUMMARY: 16 checks, 16 PASS, 0 FAIL'' (home-etch rules); ``SUMMARY: 80
checks, 80 PASS, 0 FAIL'' (interfaces) \ev{reports/electronics/evidence/T-073/};
\nolinkurl{check_gap.out} ``SUMMARY: copper PASS (4 kV rule, (a) and (b)); end caps
PARTIAL at 4 kV over the swept ASSUMPTIONS; mounting metal PASS''. These fixes were
not re-reviewed independently.

\paragraph{What the re-review did not cover.} The findings file does not list
exclusions. It is a design-file and document review: nothing in it is a hardware
measurement, and it did not see the T-073 fixes or the systems update T-074.
